\documentclass[10pt]{article}
\usepackage[letterpaper,margin=0.78in]{geometry}
\usepackage{fontspec}
\setmainfont{texgyreschola}[
  Extension=.otf,
  UprightFont=*-regular,
  BoldFont=*-bold,
  ItalicFont=*-italic,
  BoldItalicFont=*-bolditalic
]
\usepackage{microtype}
\setlength{\emergencystretch}{3em}
\hfuzz=16pt
\usepackage{ragged2e}
\usepackage{array,longtable,tabularx}
\usepackage{multirow}
\usepackage{enumitem}
\usepackage{xcolor}
\definecolor{cvlinkblue}{HTML}{00AEEF}
\usepackage[colorlinks=true,urlcolor=cvlinkblue,linkcolor=cvlinkblue,citecolor=cvlinkblue]{hyperref}
\urlstyle{same}
\setlength{\parindent}{0pt}
\setlength{\parskip}{0pt}
\setlength{\LTpre}{0pt}
\setlength{\LTpost}{0pt}
\renewcommand{\arraystretch}{1.05}
% Temporary column guides: change to \cvshowcolumnsfalse to hide them.
\newif\ifcvshowcolumns
\cvshowcolumnsfalse
\newcommand{\cvcolumnguide}{%
  \ifcvshowcolumns{\color{red!65}\vrule width 0.3pt\kern-0.3pt}\fi}
\newcolumntype{L}{>{\raggedright\arraybackslash\bfseries}p{0.17\textwidth}}
\newcolumntype{R}{>{\raggedright\arraybackslash}p{\dimexpr0.65\textwidth-2\tabcolsep\relax}}
\newcolumntype{D}{>{\raggedleft\arraybackslash\itshape}p{0.12\textwidth}}
% Rows without a separate date use the full content width.
\newcommand{\cvwide}[1]{\multicolumn{2}{>{\justifying\arraybackslash}p{0.77\textwidth}@{}}{#1}}
% Institution or venue, followed by the location when provided.
% Place this row immediately after the entry title and before its description.
\newcommand{\cvlocation}[1]{\cvwide{\noindent\textit{#1}}}
\newcommand{\cvblank}{\rule{0pt}{0.65em}}
% Indented award comments; continuation lines align with the comment text.
\newcommand{\awardcomment}[1]{%
  \cvwide{\noindent\hspace*{1em}\parbox[t]{\dimexpr\linewidth-1em\relax}{%
    \justifying\hangindent=1em\hangafter=1\noindent
    \makebox[1em][l]{--}#1}}}

\begin{document}
\thispagestyle{empty}
\begin{center}
{\LARGE\bfseries Nara Ávila}\\[1.1em]
{\large Master's Student}\\
Department of Mathematical Physics\\
Physics Institute, University of São Paulo (USP)\\
\emph{E-mail:} \href{mailto:nara.moraes@usp.br}{nara.moraes@usp.br}\\
\emph{Website:} \url{https://fma.if.usp.br/~avilamrs}
\end{center}
\vspace{1.0em}
\hrule
\vspace{1.2em}
\small
\begin{longtable}{@{}L!{\cvcolumnguide}R!{\cvcolumnguide}D@{}}

Education 
 & \textbf{M.Sc.} in Nuclear Physics \href{https://fma.if.usp.br/~avilamrs}{[link]} & 2025 -- Present \\*
 & \cvlocation{Physics Institute (IF), University of São Paulo (USP)} \\
 & \cvwide{\cvblank} \\[-0.35em]
& \textbf{B.Sc.} in Computer Engineering
  \href{https://naavilam.github.io/diploma/Computer-Engineer-Bachelor.pdf}{[link]}
  & 2006 -- 2011 \\*
& \cvlocation{São Carlos School of Engineering (EESC), University of São Paulo (USP)} \\*

& \cvwide{The interdisciplinary curriculum encompassed four major areas of
computer engineering:} \\

& \cvwide{\textbf{Computer Science \& Software Engineering:}
Algorithms and data structures, compilers, object-oriented and distributed
programming, software engineering, UML, design patterns, MVC, software testing,
relational databases and SQL, linear optimization, and performance evaluation.} \\

& \cvwide{\textbf{Computer Architecture:}
Digital logic, microprocessors, pipelined RISC/CISC architectures, parallel
computing, SIMD/MIMD systems, and shared- and distributed-memory organization.
Hardware design included VHDL, Verilog, programmable logic, and ASIC synthesis
and layout using LeonardoSpectrum and Mentor tools.} \\

& \cvwide{\textbf{Physics \& Electronics:}
Classical and quantum mechanics, electromagnetism, circuit analysis,
semiconductor devices, analog and digital electronics, CMOS microelectronics,
digital signal processing, Fourier and Z transforms, digital filters, and
microprocessor-based control systems.} \\

& \cvwide{\textbf{Communications \& Networks:}
Communication theory, modulation, antennas, optical communications, TCP/IP,
routing, TCP/UDP, wireless networks, and quality of service. Practical work
included Java video streaming with RTSP/RTP, IEEE 802.11 analysis with Wireshark,
and Android networking with Bluetooth and Wi-Fi.} \\

& \cvwide{\cvblank} \\[-0.35em]
\multirow{2}{\hsize}{Academic Positions}
& \textbf{Research Collaborator}
  \href{https://naavilam.github.io/reports/RCGI-Report.pdf}{[link]}
  & 2023 -- 2025 \\*
& \cvlocation{Department of Materials Physics and Mechanics (DFMT), Institute of Physics (IF), University of São Paulo (USP)} \\*
& \cvwide{\textit{40 hours/week}} \\
& \cvwide{\textit{Project: Quantum Computing Applied to Molecular Modeling}} \\

& \cvwide{\textbf{Research Context:}
Investigation of quantum computing methods for first-principles molecular
modeling in materials and energy research. The work focused on computing
ground-state energies of molecules relevant to the energy transition, including
H$_2$ and CO$_2$, as a first step toward the use of quantum simulations in the
discovery and development of new materials and sustainable energy technologies.} \\

& \cvwide{\textbf{Quantum Simulation of Molecular Systems:}
Developed first-principles quantum simulations for molecular ground-state
energy estimation. Electronic-structure problems were formulated within the
Born--Oppenheimer approximation and second quantization, with fermionic
Hamiltonians mapped onto qubit operators using the Jordan--Wigner transformation.
Ground-state energies were estimated with VQE using UCCSD ansätze and benchmarked
against Hartree--Fock and Full Configuration Interaction (FCI) calculations.
The study examined basis-set selection, active-space approximations, shot counts,
and qubit-reduction strategies, and explored restricted Boltzmann machines (RBMs)
as alternative variational representations of molecular quantum states.} \\

& \cvwide{\textbf{Quantum Computing Methods \& Platforms:}
Implemented and evaluated gate-based quantum algorithms across Amazon Braket,
IBM Qiskit, Google Cirq, CERN Qibo, and PennyLane, considering superconducting,
trapped-ion, photonic, and neutral-atom architectures. Hardware-oriented studies
examined execution time, noise, and decoherence, together with NISQ error
mitigation using zero-noise extrapolation (ZNE) and Clifford data regression
(CDR).} \\



& \cvwide{\cvblank} \\[-0.35em]

 & \textbf{Teaching Assistant} \href{https://naavilam.github.io/agreements/SME0300-TA.pdf}{[link]} & Fall 2009 \\*
 & \cvlocation{Institute of Mathematics and Computer Sciences (ICMC), University of São Paulo (USP)} \\*
 & \cvwide{\textit{8 hours/week}} \\
 & \cvwide{\textbullet{} SME0300 Numerical Methods I} \\
 & \cvwide{\cvblank} \\[-0.35em]

 & \textbf{Teaching Assistant} \href{https://naavilam.github.io/agreements/SME0601-TA.pdf}{[link]} & Spring 2009 \\*
 & \cvlocation{Institute of Mathematics and Computer Sciences (ICMC), University of São Paulo (USP)} \\*
 & \cvwide{\textit{8 hours/week}} \\
 & \cvwide{\textbullet{} SME0601 Numerical Methods II} \\
 & \cvwide{\cvblank} \\[-0.35em]
 & \cvwide{\cvblank} \\[-0.35em]
\multirow{2}{\hsize}{Academic Awards}
 & \textbullet{} Best Presentation in Environmental Management System (SHS0436) & Fall 2009 \\*[0.25em]
 & \awardcomment{Recognized for an impactful final course presentation entitled ``Harmonic and Disharmonic Systems,'' due to its combination of visual simplicity and strong conceptual content. Using images alone, the presentation explored in detail the diversity, complexity, and beauty of harmonic systems in nature---including coral ecosystems and harmonious ecological relationships such as mutualism, commensalism, and epiphytism---in contrast with human systems that disrupt environmental balance, illustrated through examples of waste accumulation, pollution, environmental degradation, and pollution-related wildlife mortality, alongside other disharmonic relationships also found in nature, such as parasitism and disease. Received the maximum final grade (10.0/10.0).} \\\noalign{\vskip 0.65em}
& \textbullet{} 9th place -- ACM-ICPC / SBC Brazilian Programming Contest,
First Phase
\href{https://maratona.sbc.org.br/hist/2007/primeira-fase/}{[link]}
& Sep 2007 \\*[0.25em]
& \cvlocation{Campinas} \\*
& \awardcomment{Competed as a member of team ``Greve de Acertos?'' in the
first phase of the Brazilian Programming Contest
(Maratona de Programação), the Brazilian regional competition of the
ACM International Collegiate Programming Contest (ICPC).} \\
\noalign{\vskip 0.65em}

& \textbullet{} 6th place in the IPhO / SBF Brazilian Physics Olympiad
\href{https://naavilam.github.io/certificates/obf-2003.pdf}{[link]}
& Oct 2003 \\*[0.25em]
& \awardcomment{National physics olympiad organized by the Brazilian Physical Society (SBF) and part of Brazil's selection process for the International Physics Olympiad (IPhO).} \\
\noalign{\vskip 0.65em}

& \textbullet{} 18th place in the OBM / XI OMEG
\href{https://naavilam.github.io/certificates/omeg-2002.pdf}{[link]}
& Sep 2002 \\*[0.25em]
& \awardcomment{Mathematics olympiad competition associated with the Brazilian Mathematical Olympiad (OBM), Brazil's national mathematics olympiad and part of the country's selection pathway for the International Mathematical Olympiad (IMO).} \\
\noalign{\vskip 0.65em}

& \cvwide{\cvblank} \\[-0.35em]
& \cvwide{\cvblank} \\[-0.35em]
\pagebreak[2]

Grants & \textcolor{gray!60}{\textbf{AWS Cloud Credit for Research}}
\href{https://pages.awscloud.com/aws-cloud-credit-for-research.html}{[link]}
& \textcolor{gray!60}{Sep 2026} \\*
& \cvwide{\textcolor{gray!60}{Research Area: Quantum Computing}} \\*
& \cvwide{\textcolor{gray!60}{Requested Amount: US\$5,000.00 in AWS Braket credits}} \\*
& \cvwide{\textcolor{gray!60}{Status: Submitted September 17, 2026}} \\
& \cvwide{\cvblank} \\[-0.35em]
& \textbf{EAP5054 Entrepreneurial Training for Scientists II (PRPG/USP)} \href{https://prpg.usp.br/pt-br/pae/oficina-empreendedorismo}{[link]} & Spring 2025 \\
& \cvwide{Project Title: StrategiQ --- AI-Based SaaS for Quantum Computing Governance Consulting} \\
& \cvwide{Role: Founder and Project Lead of a three-member team conducting market research and developing the business plan for StrategiQ developed through the Entrepreneurial Training for Scientists I (DPG5011)} \\
& \cvwide{Funding Agency: Office of the Provost for Graduate Studies, University of São Paulo (PRPG/USP)} \\
& \cvwide{Total Amount: R\$3,720.00} \\
& \cvwide{\cvblank} \\[-0.35em]
& \textbf{132734/2025-7} \href{https://naavilam.github.io/grants/132734-2025-7.pdf}{[link]} & 2025--2027 \\*
& \cvlocation{University of São Paulo (USP)} \\*
& \cvwide{Program: Master's Scholarship (GM)} \\
& \cvwide{Funding Agency: National Council for Scientific and Technological Development (CNPq)} \\
& \cvwide{Total Amount: R\$50,400.00} \\
& \cvwide{\cvblank} \\[-0.35em]

& \textbf{409642/2022-3} \href{https://naavilam.github.io/grants/409642-2022-3.pdf}{[link]} & 2023--2024 \\
& \cvwide{Project Title: Quantum Computing Applied to Molecular Modeling} \\
& \cvwide{Program: Technological Development in ICTs (DTC-F)} \\
& \cvwide{Funding Agency: National Council for Scientific and Technological Development (CNPq)} \\
& \cvwide{Scholarship Process: 301681/2023-6} \\
& \cvwide{Total Amount: R\$36,300.00} \\
& \cvwide{\cvblank} \\[-0.35em]

\pagebreak[2]

Software
& 1. MORAES, N. A.; LUZUM, M. W. \textit{hic-bayes}.
\href{https://fma.if.usp.br/~avilamrs/index.html\#report-2026-06}{[Sec. 4.1]}
& 2026 \\
& \cvwide{Python software for Bayesian parameter estimation and Gaussian-process emulation in relativistic heavy-ioncollision studies. Implements PCA-based dimensionality reduction, Gaussian-process emulators, multivariate likelihoods, MCMC sampling, posterior analysis, and multi-fidelity workflows.} \\
& \cvwide{\cvblank} \\[-0.4em]
& 2. MORAES, N. A. \textit{Zotero Mirroring Plugin}.
\href{https://naavilam.github.io/_posts/2026-02-02-Zotero-Plugin.html}{[link]}
& 2026 \\
& \cvwide{Zotero plugin for mirroring the hierarchical structure of a Zotero library onto the local filesystem.} \\
& \cvwide{\cvblank} \\[-0.4em]
& 3. MORAES, N. A. \textit{Git Release and Milestone CLI}.
\href{https://naavilam.github.io/_posts/2025-12-12-GitHub-Version-Control.html}{[link]}
& 2026 \\
& \cvwide{Command-line tool for automating Git releases, milestones, and software development cycles.} \\
& \cvwide{\cvblank} \\[-0.4em]
& 4. MORAES, N. A. \textit{Quantum Battleship}.
\href{https://naavilam.github.io/_posts/2025-04-11-Quantum-Battleship.html}{[link]}
& 2025 \\
& \cvwide{Quantum computing game that allows users to play Battleship using real quantum computers.} \\
& \cvwide{\cvblank} \\[-0.4em]
& 5. MORAES, N. A. \textit{Molecular Symmetry}.
\href{https://naavilam.github.io/_posts/2025-06-16-Molecular-Symmetry.html}{[link]}
& 2025 \\
& \cvwide{Web application for parsing molecular geometries from XYZ files and generating molecular symmetry reports.} \\


& \cvwide{\cvblank} \\[-0.35em]
& \cvwide{\cvblank} \\[-0.35em]
\pagebreak[2]
\multirow{2}{\hsize}{Leadership Experience} & Founder \& Team Lead \href{https://quantumera.ai/strategiq/index.html}{[link]} & Fall 2025 \\*
 & \cvlocation{University of São Paulo (USP)} \\*
 & \cvwide{DPG5011 -- Entrepreneurship} \\
 & \cvwide{\textbullet{} Founded and led a six-member interdisciplinary team in the development of an} \\
 & \cvwide{innovative startup focused on quantum computing.} \\
 & \cvwide{\textbullet{} Led the team through the development of the startup concept, including problem} \\
 & \cvwide{definition, business strategy, and the application of quantum technologies to} \\
 & \cvwide{industry-oriented solutions.} \\
 & \cvwide{\cvblank} \\[-0.35em]

 & Founder \& Team Lead \href{https://quantumera.ai/strategiq/index.html}{[link]} & Spring 2025 \\*
  & \cvlocation{University of São Paulo (USP)} \\*
 & \cvwide{EAP5054 -- Entrepreneurship} \\
 & \cvwide{\textbullet{} Founded and led a three-member interdisciplinary team in the development of an} \\
 & \cvwide{innovative startup focused on quantum computing.} \\
 & \cvwide{\textbullet{} Led the team through the development of the startup concept, including problem} \\
 & \cvwide{definition, business strategy, and the application of quantum technologies to} \\
 & \cvwide{industry-oriented solutions.} \\
 & \cvwide{\cvblank} \\[-0.35em]

 & Undergraduate Research Mentor & 2025 \\*
 & \cvlocation{Institute of Physics, University of São Paulo (IFUSP)} \\*
 & \cvwide{\textbullet{} Mentored a group of five undergraduate students in scientific initiation} \\
 & \cvwide{projects in quantum computing, providing guidance throughout their research activities.} \\
 & \cvwide{\textbullet{} Supervised the development of individual research projects presented by the} \\
 & \cvwide{students at the Undergraduate Research Symposium.} \\
& \cvwide{\cvblank} \\[-0.35em]
& \cvwide{\cvblank} \\[-0.35em]

\multirow{2}{\hsize}{Applications \& Proposals}

& \textcolor{gray}{Quantum Computing and Sustainability:
Hackathon and Workshop}
\href{https://www.rioquantumhack.com.br/en}{[link]}
& \textcolor{gray}{Nov--Dec 2026} \\*[0.25em]
& \cvlocation{\textcolor{gray}{Brazilian Center for Research in Physics (CBPF) -- Rio de Janeiro}} \\*
& \awardcomment{\textcolor{gray}{Application accepted; participation in the Quantum
Computing and Sustainability Hackathon confirmed.}} \\
\noalign{\vskip 1.25em}

& Application to MBA Program \href{https://www.hbs.edu/mba}{[link]}
& Jan 2023 \\*[0.25em]
& \cvlocation{Harvard Business School (HBS)} \\*
& \awardcomment{Submitted application for admission to the MBA Program.} \\
\noalign{\vskip 0.65em}
 & \cvwide{\cvblank} \\[-0.35em]
 & \cvwide{\cvblank} \\[-0.35em]


\multirow{2}{\hsize}{Professional Diplomas} & Postgraduate Specialization in Financial and Capital Markets \href{https://naavilam.github.io/diploma/Finance-and-Capital-Markets-Specialist.pdf}{[link]} & 2021 -- 2022 \\*
 & \cvlocation{Mackenzie Presbyterian University} \\*
 & \cvwide{\textbullet{} Quantitative methods, corporate finance, financing, and financial products and markets} \\
 & \cvwide{\cvblank} \\[-0.35em]
 & \cvwide{\cvblank} \\[-0.35em]


\multirow{2}{\hsize}{Professional Experience} & Software Engineering Specialist & 2020 -- 2022 \\*
 & \cvlocation{Open Co \href{https://open-co.com}{[link]}} \\*
 & \cvwide{\textbullet{} Developed and optimized backend services for a digital banking platform,} \\
 & \cvwide{including the first installment payment solution via PIX.} \\
 & \cvwide{\textbullet{} Proposed and prototyped the migration of a credit-analysis workload from} \\
 & \cvwide{Kubernetes (EKS) to AWS Lambda, reducing response time and operational costs by} \\
 & \cvwide{approximately 90\%.} \\
 & \cvwide{\cvblank} \\[-0.35em]

 & Senior Software Engineer & 2019 -- 2021 \\*
 & \cvlocation{TIDEXA \href{https://www.youtube.com/user/tidexa}{[link]}} \\*
 & \cvwide{\textbullet{} Contributed to the migration of a payroll system from Oracle Forms/PL--SQL} \\
 & \cvwide{to a Java-based platform.} \\
 & \cvwide{\textbullet{} Led technical decisions and development activities within the engineering team.} \\
 & \cvwide{\cvblank} \\[-0.35em]

 & Competitive Public-Sector Examination Preparation & 2016 -- 2018 \\
 & \cvwide{Independent, full-time study} \\
 & \cvwide{\textbullet{} Multidisciplinary preparation for competitive Brazilian public-sector} \\
 & \cvwide{examinations, including Law, Accounting, Economics, Public Finance, Public} \\
 & \cvwide{Administration, and Business Administration.} \\
 & \cvwide{\cvblank} \\[-0.35em]

 & Mid-Level Software Engineer & 2014 -- 2016 \\*
 & \cvlocation{Itautec -- Grupo ITAUSA \href{https://pt.wikipedia.org/wiki/Itautec}{[link]}} \\*
 & \cvwide{\textbullet{} Redesigned a business-critical software application using software design} \\
 & \cvwide{patterns and engineering best practices, improving usability, reliability, and} \\
 & \cvwide{maintainability. The application was deployed across more than 5,000 branches of a major} \\
 & \cvwide{Brazilian bank.} \\
 & \cvwide{\cvblank} \\[-0.35em]

 & Junior Software Engineer & 2012 -- 2014 \\*
 & \cvlocation{Itautec -- Grupo ITAUSA} \\*
  & \cvwide{\textbullet{} Developed and maintained software applications for retail industry.} \\

 & \cvwide{\cvblank} \\[-0.35em]
& \cvwide{\cvblank} \\[-0.35em]


\multirow{2}{\hsize}{Talks \& Presentations}
& \cvwide{\textbf{Conference Presentations}} \\[0.3em]

& \cvwide{1. ``The Effect of Correlated Experimental Uncertainties on Bayesian Inference of Quark--Gluon Plasma Properties.'' \href{https://fma.if.usp.br/~avilamrs/index.html\#phenoexp}{[link]}} \\*
& \cvlocation{PHENOEXP Workshop} \\*
& \cvwide{Poster presentation, August 2026.} \\[0.4em]

& \cvwide{2. ``Study of the Effect of Correlated Experimental Uncertainty on Bayesian Inference of Quark--Gluon Plasma Properties.'' \href{https://fma.if.usp.br/~avilamrs/index.html\#retinha-xxxiii}{[link]}} \\*
& \cvlocation{XXXIII Reunião de Trabalho sobre Interações Hadrônicas (RETINHA XXXIII)} \\*
& \cvwide{Oral presentation, June 2026.} \\[0.4em]

& \cvwide{3. ``Risk Management and IT Governance in the Adoption of Quantum Computing: A COBIT 5-Based Conceptual Model.'' \href{https://naavilam.github.io/paper/quantum_it_governance_en.pdf}{[link]}} \\*
& \cvlocation{5th Congress on Risk Management and ERM -- FEA--USP} \\*
& \cvwide{Oral presentation, December 2025.} \\[0.4em]

& \cvwide{4. ``Exploring the Potential of Quantum Computing in Materials Science.'' \href{https://naavilam.github.io/certificates/20250518_EOSBF.pdf}{[link]}} \\*
& \cvlocation{Autumn Meeting of the Brazilian Physical Society} \\*
& \cvwide{Oral presentation (co-author), May 2025.} \\[0.4em]
& \cvwide{\cvblank} \\[-0.35em]
& \cvwide{\textbf{Industry \& Professional Talks}} \\[0.3em]

& \cvwide{1. ``Unveiling Quantum Computing: Practical Applications and Future Possibilities.'' \href{https://naavilam.github.io/assets/img/eventos/evento4.png}{[link]}} \\*
& \cvlocation{Rio Innovation Week -- AI and Meta World} \\*
& \cvwide{Talk, 2024.} \\[0.4em]

& \cvwide{2. ``Quantum Computing. Hands On: Grover Search Algorithm.''} \\*
& \cvlocation{Open Co -- Tech Day} \\*
& \cvwide{Talk, 2022. \textit{[Internal material]}} \\[0.4em]

& \cvwide{3. ``Shift Your Programming Paradigm to a Quantum Mindset.'' \href{https://naavilam.github.io/presentations/20220822_TDC_Business.pdf}{[link]}} \\*
& \cvlocation{The Developer's Conference (TDC Business)} \\*
& \cvwide{Talk, 2022.} \\[0.4em]

& \cvwide{4. ``Fundamentals of Quantum Computing.''} \\*
& \cvlocation{Open Co -- Tech Day} \\*
& \cvwide{Talk, 2021. \textit{[Internal material]}} \\[0.4em]

& \cvwide{5. ``Technical Software Documentation.''} \\*
& \cvlocation{ITAUTEC} \\*
& \cvwide{Seminar, 2014. \textit{[Internal material]}} \\[0.4em]
& \cvwide{\cvblank} \\[-0.35em]
& \cvwide{\cvblank} \\[-0.35em]

\multirow{2}{\hsize}{Early Education Highlights}

& School Environmental Research Project -- Municipal Waste Management & 1999 \\*
& \cvlocation{Colégio Galáxia -- Goiânia, Goiás, Brazil} \\*
& \cvwide{\textbullet{} Conducted a school research project on municipal solid waste management, including} \\
& \cvwide{an interview with Pedro Wilson Guimarães, then Federal Deputy for the State of Goiás.} \\
& \cvwide{\textbullet{} Investigated waste-management challenges and technologies, including landfill} \\
& \cvwide{biogas recovery and international approaches to waste incineration.} \\
& \cvwide{\textbullet{} Wilson later served as Mayor of Goiânia (2001--2004), where his administration} \\
& \cvwide{established municipal selective waste collection through legislation enacted in 2002.} \\
& \cvwide{\cvblank} \\[-0.35em]
& \cvwide{\cvblank} \\[-0.35em]

& School Environmental Field Visit -- Goiânia Municipal Landfill \href{https://aterrogoiania.comurg.com.br}{[link]} & 1998 \\*
& \cvlocation{Escola do Parque \href{https://www.escolaparquegoiania.com.br}{[link]} -- Goiânia, Goiás, Brazil} \\*
& \cvwide{\textbullet{} Participated in a field visit to the Goiânia Municipal Landfill, observing} \\
& \cvwide{municipal solid waste disposal operations and environmental-control systems.} \\
& \cvwide{\textbullet{} Visited the landfill-gas collection and flaring system, where biogas was captured} \\
& \cvwide{and burned rather than recovered for energy production at the time.} \\
& \cvwide{\textbullet{} Observed the leachate treatment system used to treat landfill effluent before} \\
& \cvwide{the treated water was discharged back into the local river system.} \\
& \cvwide{\cvblank} \\[-0.35em]

& ``Aprender a Pensar'' -- High Abilities/Gifted Education Extension Program & 1996 \\*
& \cvlocation{Pontifical Catholic University of Goiás (PUC Goiás) -- Goiânia, Goiás, Brazil} \\*
& \cvwide{\textbullet{} Selected for participation in a university extension program for children identified} \\
& \cvwide{with high abilities/giftedness through psychometric assessment of non-verbal abstract} \\
& \cvwide{reasoning, including visual pattern recognition and inductive reasoning tasks.} \\
& \cvwide{\textbullet{} Participated in enrichment activities designed to develop higher-order reasoning,} \\
& \cvwide{problem-solving, creativity, and metacognitive skills.} \\
& \cvwide{\cvblank} \\[-0.35em]
& \cvwide{\cvblank} \\[-0.35em]

& Logosophical Education -- First Grade & 1995 \\*
& \cvlocation{Colégio Logosófico González Pecotche -- Goiânia, Goiás, Brazil} \\*
& \cvwide{\textbullet{} Completed first grade at a school whose educational approach was grounded in} \\
& \cvwide{Logosophy, a system of thought developed by Carlos Bernardo González Pecotche that} \\
& \cvwide{emphasizes conscious self-knowledge and the development of individual capacities.} \\
& \cvwide{\textbullet{} The educational experience placed particular emphasis on activities fostering} \\
& \cvwide{self-awareness, reflection, creativity, and personal expression.} \\
& \cvwide{\cvblank} \\[-0.35em]
& \cvwide{\cvblank} \\[-0.35em]

Certificates
& \textcolor{gray!60}{\textbf{HSK 3} -- Chinese Proficiency Exam Scheduled}
\href{https://institutoconfucio.com.br/exame-de-proficiencia/}{[link]}
& \textcolor{gray!60}{18 Oct 2026} \\*
& \cvlocation{\textcolor{gray!60}{Instituto Confúcio -- São Paulo, SP}} \\
& \cvwide{\cvblank} \\[-0.35em]
& \cvwide{\cvblank} \\[-0.35em]

\multirow{2}{\hsize}{Extracurricular Activities} & 1st Place -- Brazil Interstate Taekwon-Do ITF Competition \href{http://www.fbt.org.br/resultados_det.asp?id=54}{[link]} & Aug 2023 \\*
 & \cvlocation{Búzios, Rio de Janeiro, Brazil} \\*
 & \cvwide{\textbullet{} Two Gold medals: Sparring and Patterns (Tul).} \\
 & \cvwide{\cvblank} \\[-0.35em]

  & 1st Place -- Children's Drawing Contest \href{https://thy-catalogue.art/galerie}{[Art Portifolio]} & Oct 2000 \\*
 & \cvlocation{Claro (telecommunications company) -- Brazil} \\*
 & \cvwide{\textbullet{} First-place award in a Children's Day drawing competition organized by} \\
 & \cvwide{Claro for employees' children.} \\

\end{longtable}
\vfill
{\footnotesize\itshape\textsuperscript{\tiny *} Information shown in gray refers to future or pending items.\par}
\end{document}
